\section{Conclusion}
\label{sec:conclusion}

\subsection{Summary}
\label{subsec:summary}
This work presents an architecture drawing on approaches by Kim et al. (2021)~\cite{Kim2021} and Kolter et al. (2020)~\cite{Kolter2020}
to optimize the parameters of multiple SABR models using a neural network with implicit deep equilibrium layers as the candidate set selection mechanism.
Our approach was successful in enforcing the no-arbitrage conditions in the resulting surfaces by solving for a fixed-point problem between the observed data and generated IV values.
The time complexity of the trained model vastly outperformed the Monte Carlo method, with a speedup of 6300x, and was able to generalize well to test data,
as shown in \hyperref[tab:performance]{Table~\ref*{tab:performance}} and \hyperref[tab:training]{Table~\ref*{tab:training}}, respectively.
Additionally, the neural network approach supports fine-tuning to incorporate new daily observations;
a significant advantage over the Monte Carlo process, which would require estimation from scratch using the new dataset.

\subsection{Implications and Future Work}
\label{subsec:implications}
The main focus of this work was to merge the approach using the transformer and the implicit layers to
ensure no-arbitrage conditions in the SABR model by solving for a fixed-point problem.
The approach of using neural networks, even though significantly faster than Monte Carlo methods, could be further optimized for the context of high-frequency intraday trading.
Adapting the model to accept exogenous market factors besides the observed data, such as macroeconomic indicators, market indexes, and other financial instruments,
could allow the model to better adapt to changing market regimes and is also a promising extension in future research.
